\section{Limitations}

\paragraph{Evaluation scope.}
The current manuscript reports current ReCEG values in three Amazon
target-domain tables. The tables should not be interpreted as seed-robust
mean-and-standard-deviation estimates or as statistical-significance evidence.

\paragraph{Reference and benchmark compatibility.}
Strict longest-reference reconstruction and full D4C paper-compatible
benchmark reproduction are outside the present paper boundary. The current
tables are organized for compact method presentation, with baseline values
adapted from D4C-style reporting and ReCEG values shown as current bounded
evidence.

\paragraph{Broader baselines.}
Modern external baselines such as CIER, MAPLE, ELIXIR, Prism, and XRec require
same-setting adaptation before they can be used as result-table comparisons.
They are therefore discussed as comparison scope rather than as current result
claims.

\paragraph{Ablation scope.}
The current ablation table focuses on AM-CDs and on the two controls requested
for the main paper story: Semantic Evidence modeling and Reliability Guidance.
The manuscript therefore treats the ablation as table-supported evidence for
those controls, while avoiding broader component-isolation claims that would
require additional target-domain ablations.

\paragraph{Figure and presentation polish.}
The core method figures are now vector/TikZ diagrams intended for manuscript
review. Their visual density, placement, and AAAI camera-ready styling may
still need refinement after Chat review.
